% संपूर्ण मराठी ग्रंथातील प्रकरण 68: विगमन
% हा संपादनयोग्य स्रोतखंड आहे. संकलनासाठी संपूर्ण स्रोतसंग्रहातील
% अक्षररूपे, पूर्वभाग, चिन्हव्याख्या आणि आकृती-स्रोत आवश्यक आहेत.
% मूळ: Open Logic Project; मजकूर आणि रूपांतर CC BY 4.0.
% यंत्रानुवाद, दुरुस्ती आणि तपासणी: OpenAI Codex —
% GPT-5.6 Sol आणि GPT-6 Sol, दोन्ही Ultra effort.
% दुरुस्ती-पुनरावलोकन: GPT-6.1 Sol, Ultra effort.
\chapter{विगमन}\label{mth:ind:chap}









\section{प्रस्तावना}\label{mth:ind:int:sec}

विगमन ही सिद्धतेची महत्त्वाची पद्धती आहे. तर्कशास्त्र, सैद्धान्तिक संगणकशास्त्र आणि गणित यांच्या जवळजवळ सर्व शाखांत ती वेगवेगळ्या रूपांत वापरली जाते. तर्कशास्त्रातील अनेक निष्कर्ष सिद्ध करण्यासाठी ती आवश्यक आहे.

विगमनाचा अनेकदा निगमनाशी विरोध दाखवला जातो आणि विशिष्ट प्रकरणांकडून सामान्य विधानाकडे जाणारे अनुमान, असे त्याचे वर्णन केले जाते. उदाहरणार्थ, अनेक हिरवी पाचूची रत्ने आपण पाहिली आणि पाचू म्हणता येईल असे कोणतेही रत्न हिरवे नसलेले आढळले नाही, तर सर्व पाचू हिरवे असतात असा निष्कर्ष काढू शकतो. हे विगमनाधारित अनुमान आहे: ते अनेक विशिष्ट प्रकरणांकडून---हा पाचू हिरवा आहे, तो पाचू हिरवा आहे, इत्यादी---एका सामान्य विधानाकडे---सर्व पाचू हिरवे असतात---जाते. \emph{गणितीय} विगमनातही सामान्य विधान निष्कर्ष म्हणून मिळते; पण ते निरीक्षणांवरून केलेल्या या “साध्या विगमना”पेक्षा फार वेगळे असते.

स्थूलपणे सांगायचे तर, गणितातील विगमनाधारित सिद्धतेने एखाद्या प्रकारच्या सर्व गणितीय वस्तूंमध्ये विशिष्ट गुणधर्म आहे असा निष्कर्ष मिळतो. सर्वात साध्या प्रकरणात या सिद्धतेतील गणितीय वस्तू नैसर्गिक संख्या असतात. अशा वेळी सर्व नैसर्गिक संख्यांत एखादा गुणधर्म आहे हे सिद्ध करण्यासाठी विगमन वापरतो; त्यात असे दाखवतो:
\begin{enumerate}
\item $0$ मध्ये तो गुणधर्म आहे; आणि
\item एखाद्या~$k$ संख्येत तो गुणधर्म असेल, तर~$k+1$ मध्येही तो असतो.
\end{enumerate}
ज्या गणितीय वस्तूंना संख्या देता येतात त्यांच्याविषयीची सामान्य विधाने सिद्ध करण्यासाठीही नैसर्गिक संख्यांवरील विगमन अनेकदा वापरता येते. उदाहरणार्थ, प्रत्येक सांत संचात सांत संख्येने, समजा~$n$, घटक असतात. विगमन वापरून प्रत्येक~$n$ संख्येसाठी “~$n$ आकारमानाचे सर्व सांत संच \dots{} असतात” हा गुणधर्म आहे असे दाखवता आले, तर सर्व सांत संचांविषयी काहीतरी दाखवले असेल.

\emph{विगमनाधारित व्याख्येने दिलेल्या} गणितीय वस्तूंपर्यंतही विगमनाचा विस्तार करता येतो. उदाहरणार्थ, प्रथम-क्रम तर्कशास्त्रासारख्या औपचारिक भाषेतील अभिव्यक्ती विगमनाधारित व्याख्येने दिलेल्या असतात. \emph{संरचनात्मक विगमन} ही अशा सर्व अभिव्यक्तींविषयी निष्कर्ष सिद्ध करण्याची पद्धती आहे. विशेषतः तर्कशास्त्रात संरचनात्मक विगमन फार उपयुक्त आहे आणि मोठ्या प्रमाणात वापरले जाते.












\section{~$\Nat$ वरील विगमन}\label{mth:ind:inN:sec}

सर्वात साध्या रूपात विगमन ही सर्व नैसर्गिक संख्यांसाठी निष्कर्ष सिद्ध करण्याची पद्धती आहे. $0$ पासून सुरुवात करून वारंवार~$1$ वाढवत गेल्यास प्रत्येक नैसर्गिक संख्येपर्यंत अखेर पोहोचता येते, या तथ्याचा ती वापर करते. म्हणून एखादी गोष्ट प्रत्येक संख्येसाठी खरी आहे हे सिद्ध करण्यासाठी (1) ती $0$ साठी खरी आहे हे प्रस्थापित करता येते आणि (2) ती एखाद्या~$n$ संख्येसाठी खरी असेल तेव्हा पुढच्या~$n+1$ संख्येसाठीही खरी असते हे दाखवता येते. “~$n$ संख्येत $P$ हा गुणधर्म आहे” याचा संक्षेप $P(n)$ असा करू; तसेच “~$k$ संख्येत $P$ हा गुणधर्म आहे” याचा $P(k)$, इत्यादी. मग सर्व $n \in \Nat$ साठी $P(n)$ सिद्ध करणाऱ्या विगमनाधारित सिद्धतेत पुढील गोष्टी असतात:
\begin{enumerate}
\item $P(0)$ ची सिद्धता; आणि
\item कोणत्याही~$k$ साठी, $P(k)$ असेल तर $P(k+1)$ असते याची सिद्धता.
\end{enumerate}
हे पूर्ण स्पष्ट करण्यासाठी (1) आणि~(2) दोन्ही आहेत असे समजा. मग (1) वरून $P(0)$ सत्य आहे असे समजते. ~(2) देखील असल्यामुळे विशेषतः $P(0)$ असेल तर $P(0+1)$, म्हणजे $P(1)$, असते हे समजते. “कोणत्याही~$k$ साठी, $P(k)$ असेल तर $P(k+1)$” या सामान्य विधानात~$k$ च्या जागी~$0$ ठेवल्यावर हे मिळते. म्हणून विध्यात्मक अनुमानाने~$P(1)$ मिळते. पुन्हा (2) मध्ये, आता~$k$ च्या जागी $1$ घेतल्यावर, $P(1)$ असेल तर~$P(2)$ मिळते. ~$P(1)$ नुकतेच प्रस्थापित केलेले असल्यामुळे विध्यात्मक अनुमानाने~$P(2)$ मिळते. असेच पुढे चालू राहते. कोणत्याही~$n$ संख्येसाठी ही क्रिया $n$~वेळा केल्यावर अखेर~$P(n)$ मिळते. म्हणून (1) आणि~(2) मिळून कोणत्याही $n \in \Nat$ साठी~$P(n)$ प्रस्थापित करतात.

एक उदाहरण पाहू. $n$ फासे टाकल्यावर किती वेगवेगळ्या बेरजा मिळू शकतात हे शोधायचे आहे असे समजा. थोडे विचित्र वाटले, तरी $0$ फाशांपासून सुरुवात करू. फासेच नसतील, तर “टाकून” मिळणारी एकच बेरीज आहे: कोणतेही ठिपके नाहीत; म्हणजे बेरीज~$0$. म्हणून बेरजेच्या वेगवेगळ्या शक्यतांची संख्या~$1$ आहे. एकच फासा असेल, म्हणजे $n=1$, तर $1$ ते~$6$ अशी सहा मूल्ये शक्य आहेत. दोन फाशांनी $2$ ते $12$ पर्यंत कोणतीही बेरीज मिळू शकते; म्हणजे $11$ शक्यता. तीन फाशांनी $3$ ते $18$ पर्यंत कोणतीही संख्या मिळू शकते; म्हणजे $16$ वेगवेगळ्या शक्यता. $1$, $6$, $11$, $16$: एक नमुना दिसतो. कदाचित उत्तर $5n+1$ असेल? अर्थात, $5n+1$ ही कमाल शक्य संख्या आहे. कारण $n$ फाशांनी मिळणाऱ्या सर्वात लहान $n$ मूल्यापासून---सर्व फाशांवर $1$---सर्वात मोठ्या $6n$ मूल्यापर्यंत---सर्व फाशांवर $6$---फक्त $5n+1$ संख्या आहेत.

\begin{thm}
$n$ फाशांनी $n$ ते $6n$ यांमधील सर्व $5n+1$ शक्य मूल्ये मिळवता येतात.
\end{thm}

\begin{proof}
$P(n)$ हे विधान असू द्या: “$n$~फासे वापरून $n$ ते~$6n$ यांमधील कोणतीही संख्या मिळवता येते.” शून्य फाशांचे प्रकरण वर पाहिले आहे: रिक्त बेरीज शून्य असून एकच शक्यता असते. धन संख्यांसाठी विगमन वापरण्यास पुढील गोष्टी सिद्ध करू:
\begin{enumerate}
\item \emph{विगमनाचा आधार} $P(1)$: एकच फासा वापरून $1$ ते $6$ यांमधील कोणतीही संख्या मिळवता येते.
\item \emph{विगमनाची पायरी}: प्रत्येक धन नैसर्गिक संख्या $k$ साठी, $P(k)$ असेल तर~$P(k+1)$.
\end{enumerate}

(1) हे $6$ पृष्ठे असलेला फासा पाहून सिद्ध होते. त्याला सहाही पृष्ठे आहेत आणि $1$ ते~$6$ मधील प्रत्येक संख्या एका पृष्ठावर दिसते. म्हणून एक फासा वापरून $1$ ते~$6$ मधील कोणतीही संख्या मिळवता येते.

(2) सिद्ध करण्यासाठी सशर्त विधानाचे पूर्वांग, म्हणजे $P(k)$, गृहीत धरतो. या गृहीतकाला \emph{विगमन गृहीतक} म्हणतात. त्याचा वापर करून~$P(k+1)$ सिद्ध करतो. अवघड भाग असा: $k+1$ फाशांनी मिळणाऱ्या शक्य बेरजांचा विचार $k$~फाशांनी मिळणाऱ्या बेरजा आणि अतिरिक्त $k+1$ व्या फाशाचे मूल्य यांच्या आधारे करण्याचा मार्ग शोधायचा. विगमन गृहीतक वापरायचे असेल, तर असेच करावे लागते.

विगमन गृहीतकानुसार $k$~फाशांनी $k$ ते~$6k$ मधील कोणतीही संख्या मिळवता येते. ~$(k+1)$ व्या फाशावर~$1$ आला, तर एकूण बेरजेत $1$ वाढतो. म्हणून $k$~फासे टाकून नंतर $(k+1)$ व्या फाशावर~$1$ आणल्यास $k+1$ ते $6k+1$ मधील कोणतेही मूल्य मिळवता येते. काय उरले? $6k+2$ ते $6k+6$ ही मूल्ये. $k$ फाशांवर $6$ आणून आणि $(k+1)$ व्या फाशावर $2$ ते $6$ मधील संख्या आणून ही मूल्ये मिळवता येतात. या दोन्ही गोष्टी मिळून असे दाखवतात की $k+1$ फाशांनी $k+1$ ते $6(k+1)$ मधील कोणतीही संख्या मिळवता येते. म्हणजे~$P(k)$ हे विगमन गृहीतक वापरून~$P(k+1)$ सिद्ध केले आहे.
\end{proof}

अनेकदा वस्तूंच्या अनुक्रमाविषयी---संख्या, संच इत्यादी---काही सिद्ध करायचे असते आणि तो अनुक्रमच “विगमनाधारित” व्याख्येने दिलेला असतो. म्हणजे $(n+1)$ वी वस्तू $n$ व्या वस्तूच्या आधारे परिभाषित केलेली असते. अशा वेळी विगमन वापरतो. उदाहरणार्थ,~$n$ पर्यंतच्या नैसर्गिक संख्यांची बेरीज~$s_n$ अशी परिभाषित करता येते:
\begin{align*}
  s_0 & = 0\\
  s_{n+1} & = s_n + (n+1)
\end{align*}
या व्याख्येवरून असे मिळते:
\begin{align*}
  s_0 & = 0,\\
  s_1 & = s_0 + 1 && = 1,\\
  s_2 & = s_1 + 2 && = 1 + 2 = 3\\
  s_3 & = s_2 + 3 && = 1 + 2 + 3 = 6, \text{ इत्यादी.}
\end{align*}
आता विगमनाने $s_n = n(n+1)/2$ सिद्ध करता येते.

\begin{prop}
  $s_n = n(n+1)/2$.
\end{prop}

\begin{proof}
(1) $s_0 = 0\cdot(0 + 1)/2$ आणि (2) $s_k = k(k+1)/2$ असेल तर $s_{k+1} = (k+1)(k+2)/2$, हे सिद्ध करायचे आहे. (1) स्पष्ट आहे. (2) सिद्ध करण्यासाठी $s_k = k(k+1)/2$ हे विगमन गृहीतक घेऊ. त्याचा वापर करून $s_{k+1} = (k+1)(k+2)/2$ दाखवायचे आहे.

$s_{k+1}$ काय आहे? व्याख्येनुसार $s_{k+1} = s_k + (k+1)$. विगमन गृहीतकानुसार $s_k = k(k+1)/2$. आधीच्या समीकरणात ही किंमत ठेवता येते; मग अपूर्णांकांचे थोडे अंकगणित पुरते:
\begin{align*}
    s_{k+1} & = \frac{k(k+1)}{2} + (k+1) = {}\\
    & = \frac{k(k+1)}{2} + \frac{2(k+1)}{2} = {}\\
    & = \frac{k(k+1) + 2(k+1)}{2} = {}\\
    & = \frac{(k+2)(k+1)}{2}.
\end{align*}
\end{proof}

महत्त्वाचा धडा असा: विगमनाधारित व्याख्येने दिलेल्या $a_n$ अनुक्रमाविषयी काही सिद्ध करत असाल, तर विगमन हा स्वाभाविक मार्ग आहे. तसे नसले तरी---फाशांच्या शक्य बेरजांच्या उदाहरणाप्रमाणे---~$k+1$ चे प्रकरण~$k$ च्या प्रकरणाशी काही रीतीने जोडता आले, तर विगमन वापरता येते.












\section{प्रबल विगमन}\label{mth:ind:str:sec}

आधी चर्चिलेल्या विगमन तत्त्वात $P(0)$ सिद्ध करतो आणि $P(k)$ असेल तर $P(k+1)$ असते हेही सिद्ध करतो. दुसऱ्या भागात $P(k)$ सत्य आहे असे गृहीत धरून त्या गृहीतकावरून $P(k+1)$ सिद्ध करतो. धन निर्देशांकासाठी याऐवजी $P(k-1)$ गृहीत धरून $P(k)$ सिद्ध करता येते. महत्त्वाचे असे की कोणत्याही संख्येपासून तिच्या उत्तरवर्ती संख्येपर्यंतचे हे अनुमान करता आले पाहिजे: एखाद्या संख्येच्या पूर्ववर्ती संख्येसाठी विधान सत्य आहे असे गृहीत धरून तिच्यासाठी ते सिद्ध करता आले पाहिजे.

विगमन तत्त्वाचा आणखी एक प्रकार आहे. त्यात केवळ $k$ च्या पूर्ववर्ती $k-1$ साठी विधान खरे आहे असे न मानता,~$k$ पेक्षा लहान असलेल्या सर्व नैसर्गिक संख्यांसाठी ते खरे आहे असे मानतो आणि त्यावरून~$k$ साठी ते प्रस्थापित करतो. यानेही सर्व~$n \in \Nat$ साठी $P(n)$ मिळते. कारण $P(0)$ प्रस्थापित केल्यावर~$1$ पेक्षा लहान सर्व नैसर्गिक संख्यांसाठी $P$ खरे आहे असे प्रस्थापित झाले आहे. प्रत्येक $l<k$ साठी $P(l)$ असेल तर $P(k)$ असते, हे माहीत असल्यास विशेषतः $k=1$ साठीही ते माहीत असते. म्हणून $P(1)$ असा निष्कर्ष काढता येतो. यामुळे $P(0)$ आणि $P(1)$ सिद्ध झाले; म्हणजे सर्व $l<2$ साठी $P(l)$. तसेच प्रत्येक $l<2$ साठी $P(l)$ असेल तर $P(2)$ असते, हे सशर्त विधानही उपलब्ध असल्यामुळे~$P(2)$ निष्कर्ष म्हणून मिळते; आणि असेच पुढे.

खरे तर “प्रत्येक $k$ साठी, प्रत्येक $l<k$ साठी $P(l)$ असेल तर $P(k)$ असते” हे सामान्य सशर्त विधान प्रस्थापित करता आले, तर $P(0)$ वेगळे प्रस्थापित करण्याची गरज नसते; ते यावरूनच मिळते. “प्रत्येक $l<k$ साठी $P(l)$” सारखे सामान्य विधान कोणतेही $l<k$ नसतील तेव्हा सत्य असते, हे आठवा. हे रिक्त संख्यापनाचे प्रकरण आहे: कोणतेही $A$ नसतील, तर “सर्व $A$ हे $B$ आहेत” हे सत्य असते; कोणताही $x$ हा~$\varphi(x)$ पूर्ण करत नसेल, तर $\lforall[x][(\varphi(x) \lif \psi(x))]$ सत्य असते. या प्रकरणात औपचारिक रूप “$\lforall[l][(l < k \lif P(l))]$” असे असेल; आणि कोणतेही $l < k$ नसतील, तर ते सत्य असते. $k=0$ असेल तेव्हा नेमके हेच घडते: कोणतीही नैसर्गिक संख्या $l<0$ नाही. म्हणून $P$ काहीही असले, तरी “प्रत्येक $l<0$ साठी $P(l)$” सत्य असते. त्यामुळे “प्रत्येक $l<k$ साठी $P(l)$ असेल तर $P(k)$ असते” याची सिद्धता आपोआप~$P(0)$ प्रस्थापित करते.

~$k$ साठी विधान प्रस्थापित करताना केवळ $k-1$ साठीच्या विधानावर अवलंबून राहता येत नसेल, पण एक किंवा अधिक $l<k$ साठी ते सत्य असल्याचे गृहीतक आवश्यक असेल, तर हा प्रकार उपयुक्त ठरतो.












\section{विगमनाधारित व्याख्या}\label{mth:ind:idf:sec}

तर्कशास्त्रात अनेकदा वस्तूंचे प्रकार \emph{विगमनाधारित रीतीने} परिभाषित करतो. म्हणजे त्या प्रकारची वस्तू कशाला म्हणायचे याचे नियम देतो; त्यांत त्या प्रकारच्या जुन्या वस्तूंपासून नव्या वस्तू कशा मिळवायच्या हे सांगितलेले असते. उदाहरणार्थ, भाषेतील पदे आणि सूत्रे यांसारखे चिन्हांचे विशिष्ट अनुक्रम अनेकदा विगमनाने परिभाषित करतो. एक साधे उदाहरण पाहू: $\mathrm{a}$, $\mathrm{b}$, $\mathrm{c}$, $\mathrm{d}$ ही अक्षरे,~$\circ$ हे चिन्ह आणि $[$, $]$ हे कंस यांपासून बनलेल्या चिन्हमाला घ्या. उदाहरणार्थ, “$[[\mathrm{c} \circ \mathrm{d}][$”, “$[\mathrm{a}[]\circ]$”, “$\mathrm{a}$” किंवा “$[[\mathrm{a} \circ \mathrm{b}]\circ \mathrm{d}]$”. पहिल्या दोन चिन्हमालांत काहीतरी “चुकीचे” आहे असे वाटत असेल: पहिल्या चिन्हमालेत कंस मुळीच जुळत नाहीत; आणि “$\circ$” ने स्वतः अर्थपूर्ण असलेल्या अभिव्यक्ती “जोडल्या” पाहिजेत असे वाटेल. तिसरी आणि चौथी चिन्हमाला अधिक नीट दिसतात: प्रत्येक “$[$” साठी बंद करणारा “$]$” आहे---कंस असतील तर---आणि प्रत्येक $\circ$ च्या दोन्ही बाजूंना स्वतः नीट असलेल्या अभिव्यक्ती आहेत, ज्यांभोवती कंसांची जोडी आहे.

“नीटस पद” नेमके कशाला म्हणायचे हे निश्चित करायचे आहे. प्रथम, प्रत्येक स्वतंत्र अक्षर नीटस आहे. एकटे अक्षर नसलेले प्रत्येक नीटस पद “$[t \circ s]$” या रूपाचे असले पाहिजे; येथे $s$ आणि $t$ ही स्वतः नीटस असतात. उलट, $t$ आणि $s$ नीटस असतील, तर त्यांच्यामध्ये $\circ$ ठेवून आणि भोवती कंसांची जोडी घालून नवे नीटस पद बनवता येते. या क्रिया वापरून नीटस पदांचा संच \emph{परिभाषित} करता येईल. ही \emph{विगमनाधारित व्याख्या} आहे.

\begin{defn}[नीटस पदे]
\emph{नीटस पदांचा} संच पुढील विगमनाधारित व्याख्येने दिलेला आहे:
\begin{enumerate}
\item $\mathrm{a}$, $\mathrm{b}$, $\mathrm{c}$, $\mathrm{d}$ यांपैकी कोणतेही अक्षर नीटस पद आहे.
\item $s_1$ आणि $s_2$ नीटस पदे असतील, तर $[s_1 \circ s_2]$ देखील नीटस पद आहे.
\item इतर काहीही नीटस पद नाही.
\end{enumerate}
\end{defn}

या व्याख्येनुसार एखादी गोष्ट नीटस पद असते तेव्हाच आणि फक्त तेव्हा जेव्हा (1) आणि~(2) या दोन अटींनुसार सांत संख्येने पायऱ्यांत ती रचता येते. पहिल्या पायरीत केवळ स्वतंत्र अक्षरे असलेली सर्व नीटस पदे रचतो:
\[
\mathrm{a}, \mathrm{b}, \mathrm{c}, \mathrm{d}
\]
दुसऱ्या पायरीत रचलेल्या पदांना (2) लागू करतो. दोन अक्षरांच्या सर्व जोड्यांसाठी पुढील पदे मिळतील:
\[
[\mathrm{a} \circ \mathrm{a}], [\mathrm{a} \circ \mathrm{b}],
[\mathrm{b} \circ \mathrm{a}], \dots, [\mathrm{d} \circ \mathrm{d}]
\]
तिसऱ्या पायरीत आत्तापर्यंत रचलेल्या कोणत्याही दोन नीटस पदांना पुन्हा (2) लागू करतो. उदाहरणार्थ, $[\mathrm{a} \circ [\mathrm{a} \circ \mathrm{a}]]$ असे नवे नीटस पद मिळते: येथे $t$ हे पायरी~1 मधील $\mathrm{a}$ आहे आणि $s$ हे पायरी~2 मधील $[\mathrm{a} \circ \mathrm{a}]$ आहे. तसेच पायरी~2 मधील $[\mathrm{b} \circ \mathrm{c}]$ आणि $[\mathrm{d} \circ \mathrm{b}]$ या दोन पदांपासून $[[\mathrm{b} \circ \mathrm{c}] \circ [\mathrm{d} \circ \mathrm{b}]]$ रचता येते. असेच पुढे. या रीतीने न रचलेली कोणतीही गोष्ट नीटस पदांच्या संचात येऊ नये, हे खंड (3) निश्चित करतो.

“नीट वाटणारा” प्रत्येक चिन्हानुक्रम या व्याख्येनुसार नीटस आहे, हे अद्याप सिद्ध केलेले नाही, हे लक्षात घ्या. मात्र आपण रचू शकणारी प्रत्येक गोष्ट खरोखर “नीट वाटते”, हे स्पष्ट असले पाहिजे: कंस जुळतात आणि $\circ$ ने स्वतः नीटस असलेले भाग जोडलेले असतात.

विगमनाधारित व्याख्यांचे महत्त्वाचे वैशिष्ट्य असे की सर्व नीटस पदांविषयी काही सिद्ध करायचे असेल, तर कोणती प्रकरणे पाहावी लागतील हे व्याख्या सांगते. उदाहरणार्थ, $t$ हे नीटस पद आहे असे दिले असेल, तर $t$ कसे असू शकते हे व्याख्या सांगते: $t$ हे अक्षर असेल किंवा ते $[s_1 \circ s_2]$ या रूपाचे असेल, जिथे $s_1$ आणि~$s_2$ ही नीटस पदे आहेत. खंड (3) मुळे या दोनच शक्यता आहेत.

विगमनाधारित व्याख्येने दिलेल्या संचातील सर्व वस्तूंविषयी विधाने सिद्ध करताना प्रबल विगमन विशेष महत्त्वाचे ठरते. उदाहरणार्थ,~$n$ लांबीच्या प्रत्येक नीटस पदात $[$ ची संख्या~$< n/2$ आहे हे सिद्ध करायचे आहे असे समजा. याचा सर्व~$n$ विषयीचे विधान म्हणून विचार करता येतो: प्रत्येक $n$ साठी,~$n$ लांबीच्या कोणत्याही नीटस पदात $[$ ची संख्या $< n/2$ आहे.

\begin{prop}
कोणत्याही $n$ साठी,~$n$ लांबीच्या नीटस पदात $[$ ची संख्या $< n/2$ आहे.
\end{prop}

\begin{proof}
प्रबल विगमनाने हा निष्कर्ष सिद्ध करण्यासाठी पुढील सशर्त विधान सत्य आहे हे दाखवावे लागेल:
\begin{quote}
प्रत्येक $l < k$ साठी,~$l$ लांबीच्या कोणत्याही नीटस पदात $[$ ची संख्या $< l/2$ असेल, तर~$k$ लांबीच्या कोणत्याही नीटस पदात $[$ ची संख्या $< k/2$ असते.
\end{quote}
हे सशर्त विधान दाखवण्यासाठी त्याचे पूर्वांग सत्य आहे असे गृहीत धरा. म्हणजे कोणत्याही $l<k$ साठी,~$l$ लांबीच्या नीटस पदात $[$ ची संख्या $< l/2$ आहे असे गृहीत धरा. याला विगमन गृहीतक म्हणतो. ~$k$ लांबीच्या नीटस पदांसाठीही तेच सत्य आहे हे दाखवायचे आहे.

म्हणून $t$ हे~$k$ लांबीचे नीटस पद आहे असे समजा. नीटस पदांची व्याख्या विगमनाधारित असल्यामुळे दोन प्रकरणे आहेत: (1)~$t$ हे स्वतंत्र अक्षर आहे; किंवा (2)~$t$ हे $[s_1 \circ s_2]$ या रूपाचे आहे, जिथे $s_1$ आणि~$s_2$ ही नीटस पदे आहेत.
\begin{enumerate}
\item $t$ हे अक्षर आहे. मग $k = 1$ आणि $t$ मध्ये $[$ ची संख्या~$0$ आहे. $0 < 1/2$ असल्यामुळे विधान खरे आहे.
\item $t$ हे $[s_1 \circ s_2]$ या रूपाचे आहे, जिथे $s_1$ आणि~$s_2$ ही नीटस पदे आहेत. ~$s_1$ ची लांबी $l_1$ आणि~$s_2$ ची लांबी $l_2$ असू द्या. मग $t$ ची लांबी~$k$ ही $l_1+l_2+3$ आहे: $s_1$ आणि $s_2$ यांच्या लांबी अधिक $[$, $\circ$, $]$ ही तीन चिन्हे. $l_1+l_2+3$ हे नेहमी $l_1$ पेक्षा मोठे असल्यामुळे $l_1 < k$. त्याचप्रमाणे $l_2 < k$. म्हणजे $s_1$ आणि~$s_2$ या पदांना विगमन गृहीतक लागू होते: $s_1$ मध्ये $[$ ची संख्या~$m_1$ ही $< l_1/2$ आहे; आणि~$s_2$ मध्ये $[$ ची संख्या~$m_2$ ही $< l_2/2$ आहे.

$t$ मध्ये $[$ ची संख्या म्हणजे~$s_1$ मधील $[$ ची संख्या, अधिक~$s_2$ मधील $[$ ची संख्या, अधिक~$1$; म्हणजे $m_1 + m_2 + 1$. $m_1 < l_1/2$ आणि $m_2 < l_2/2$ असल्यामुळे:
\[
  m_1 + m_2 + 1 < \frac{l_1}{2} + \frac{l_2}{2} + 1 = \frac{l_1+l_2+2}{2} < \frac{l_1+l_2+3}{2} = k/2.
\]
\end{enumerate}
प्रत्येक प्रकरणात विगमन गृहीतकाच्या आधारे $t$ मध्ये $[$ ची संख्या $< k/2$ आहे हे दाखवले. प्रबल विगमनाने विधान सिद्ध होते.
\end{proof}

\begin{prob}
अतिनीटस पदांचा संच पुढील रीतीने परिभाषित करा:
\begin{enumerate}
\item $\mathrm{a}$, $\mathrm{b}$, $\mathrm{c}$, $\mathrm{d}$ यांपैकी कोणतेही अक्षर अतिनीटस पद आहे.
\item $s$ हे अतिनीटस पद असेल, तर $[s]$ देखील अतिनीटस पद आहे.
\item $s_1$ आणि $s_2$ अतिनीटस पदे असतील, तर $[s_1 \circ s_2]$ देखील अतिनीटस पद आहे.
\item इतर काहीही अतिनीटस पद नाही.
\end{enumerate}
~$n$ ही लांबी असलेल्या अतिनीटस पद~$t$ मध्ये $[$ ची संख्या $\le n/2 +1$ आहे हे दाखवा.
\end{prob}












\section{संरचनात्मक विगमन}\label{mth:ind:sti:sec}

आत्तापर्यंत सर्व नैसर्गिक संख्यांविषयी निष्कर्ष प्रस्थापित करण्यासाठी विगमन वापरले. पण त्याच्याशी संबंधित तत्त्व थेट वापरून विगमनाधारित व्याख्येने दिलेल्या संचातील सर्व घटक विषयी निष्कर्ष सिद्ध करता येतात. याला अनेकदा \emph{संरचनात्मक} विगमन म्हणतात; कारण ते त्या व्याख्येने दिलेल्या वस्तूंच्या संरचनेवर अवलंबून असते.

सहसा विगमनाधारित व्याख्येत (a) संचातील “आरंभीच्या” घटक ची यादी आणि (b) जुन्यांपासून नवे घटक निर्माण करणाऱ्या क्रियांची यादी दिलेली असते. उदाहरणार्थ, नीटस पदांच्या प्रकरणात आरंभीच्या वस्तू म्हणजे अक्षरे. एकच क्रिया आहे:
\begin{align*}
  o(s_1, s_2) = & [s_1 \circ s_2]
\end{align*}
नैसर्गिक संख्यांचा संच~$\Nat$ देखील विगमनाधारित व्याख्येने दिलेला आहे असा विचार करता येतो: आरंभीची वस्तू~$0$ आणि क्रिया म्हणजे~$x + 1$ हे उत्तरवर्ती फलन.

विगमनाधारित व्याख्येने दिलेल्या संचातील सर्व घटकांविषयी काही सिद्ध करायचे असेल---म्हणजे संचातील प्रत्येक घटक मध्ये~$P$ हा गुणधर्म आहे हे सिद्ध करायचे असेल---तर पुढील गोष्टी कराव्या लागतात:
\begin{enumerate}
\item आरंभीच्या वस्तूंमध्ये~$P$ आहे हे सिद्ध करा.
\item प्रत्येक~$o$ क्रियेसाठी, तिच्या सर्व आदानांत~$P$ असेल, तर फलितातही तो असतो हे सिद्ध करा.
\end{enumerate}
उदाहरणार्थ, सर्व नीटस पदांविषयी काही सिद्ध करण्यासाठी ते सर्व अक्षरांविषयी सत्य आहे हे सिद्ध करू; तसेच $s_1$ आणि $s_2$ यांविषयी ते स्वतंत्रपणे सत्य असेल, तर $[s_1 \circ s_2]$ विषयीही सत्य असते हे सिद्ध करू.

\begin{prop}
कोणत्याही नीटस पद~$t$ मध्ये $[$ ची संख्या $]$ च्या संख्येइतकी असते.
\end{prop}

\begin{proof}
संरचनात्मक विगमन वापरू. नीटस पदांची व्याख्या विगमनाधारित आहे: अक्षरे या आरंभीच्या वस्तू आणि जुन्या पदांपासून नवी नीटस पदे रचण्यासाठी $o$ ही क्रिया.
\begin{enumerate}
\item प्रत्येक अक्षरासाठी विधान सत्य आहे; कारण स्वतंत्र अक्षरात $[$ ची संख्या~$0$ आणि $]$ ची संख्याही~$0$ आहे.
\item $s_1$ मध्ये $[$ ची संख्या $]$ च्या संख्येइतकी आहे आणि $s_2$ साठीही तेच सत्य आहे असे समजा. $o(s_1, s_2)$ मध्ये---म्हणजे $[s_1 \circ s_2]$ मध्ये---$[$ ची संख्या ही $s_1$ आणि $s_2$ यांमधील $[$ च्या संख्यांची बेरीज अधिक एक आहे. $o(s_1, s_2)$ मध्ये $]$ ची संख्या ही $s_1$ आणि $s_2$ यांमधील $]$ च्या संख्यांची बेरीज अधिक एक आहे. म्हणून $o(s_1, s_2)$ मधील $[$ ची संख्या $o(s_1,s_2)$ मधील $]$ च्या संख्येइतकी आहे.
\end{enumerate}
\end{proof}

\begin{prob}
कोणत्याही नीटस पदाची सुरुवात~$]$ ने होत नाही, हे संरचनात्मक विगमनाने सिद्ध करा.
\end{prob}

संरचनात्मक विगमनाने आणखी एक सिद्धता देऊ. चिन्हांच्या~$t$ चिन्हमालेचा अरिक्त उचित पूर्वखंड म्हणजे अशी अरिक्त~$s$ चिन्हमाला की डावीकडून वाचताना प्रत्येक चिन्हानुसार ती $t$ शी जुळते, पण $t$ अधिक लांब असते. उदाहरणार्थ, $[a \circ {}$ हा $[a \circ b]$ चा अरिक्त उचित पूर्वखंड आहे. मात्र $[b \circ {}$ तसा नाही---दुसऱ्या चिन्हावर फरक आहे---आणि $[a \circ b]$ देखील तसा नाही---दोन्हींची लांबी समान आहे. येथे अरिक्त पूर्वखंडच घेतले आहेत: रिकाम्या पूर्वखंडात उघडणारे आणि बंद करणारे दोन्ही कंस शून्य असल्यामुळे पुढील काटेकोर असमानता त्यासाठी खरी नाही.

\begin{prop}\label{mth:ind:sti:prop:initial}
नीटस पद~$t$ च्या प्रत्येक अरिक्त उचित पूर्वखंडात $[$ ची संख्या $]$ पेक्षा जास्त असते.
\end{prop}

\begin{proof}
~$t$ वर विगमन करू:
\begin{enumerate}
\item $t$ हे स्वतंत्र अक्षर आहे: मग $t$ ला कोणतेही अरिक्त उचित पूर्वखंड नाहीत.
\item $t = [s_1 \circ s_2]$, जिथे $s_1$ आणि~$s_2$ ही नीटस पदे आहेत. $r$ हा $t$ चा अरिक्त उचित पूर्वखंड असेल, तर काही शक्यता आहेत:
\begin{enumerate}
\item $r$ हे केवळ $[$ आहे: मग $r$ मध्ये $[$ ची संख्या~$]$ पेक्षा एकाने जास्त आहे.
\item $r$ हे $[r_1$ आहे, जिथे $r_1$ हा~$s_1$ चा अरिक्त उचित पूर्वखंड आहे. $s_1$ नीटस पद असल्यामुळे विगमन गृहीतकानुसार $r_1$ मध्ये $[$ हे $]$ पेक्षा जास्त आहेत; आणि $[r_1$ साठीही तेच खरे आहे.
\item $r$ हे $[s_1$ किंवा $[s_1 \circ {}$ आहे. आधीच्या निष्कर्षानुसार~$s_1$ मध्ये $[$ आणि $]$ यांच्या संख्या समान आहेत. म्हणून $[s_1$ किंवा $[s_1 \circ {}$ मध्ये $[$ ची संख्या~$]$ पेक्षा एकाने जास्त आहे.
\item $r$ हे $[s_1 \circ r_2$ आहे, जिथे $r_2$ हा~$s_2$ चा अरिक्त उचित पूर्वखंड आहे. विगमन गृहीतकानुसार $r_2$ मध्ये $[$ हे $]$ पेक्षा जास्त आहेत. आधीच्या निष्कर्षानुसार~$s_1$ मध्ये $[$ आणि $]$ यांच्या संख्या समान आहेत. म्हणून $[s_1 \circ r_2$ मध्ये $[$ ची संख्या~$]$ पेक्षा मोठी आहे.
\item $r$ हे $[s_1 \circ s_2$ आहे. आधीच्या निष्कर्षानुसार $s_1$ मध्ये $[$ आणि $]$ यांच्या संख्या समान आहेत आणि~$s_2$ मध्येही त्या समान आहेत. म्हणून $[s_1 \circ s_2$ मध्ये $[$ ची संख्या~$]$ पेक्षा एकाने जास्त आहे.
\end{enumerate}
\end{enumerate}
\end{proof}












\section{संबंध आणि फलने}\label{mth:ind:rel:sec}

नैसर्गिक संख्या किंवा नीटस पदे यांसारख्या वस्तूंचा संच विगमनाधारित रीतीने परिभाषित केल्यावर, त्या वस्तूंवरील \emph{संबंध}ही विगमनाने परिभाषित करता येतात. उदाहरणार्थ, एका नीटस पद~$t_2$ मध्ये दुसरे नीटस पद~$t_1$ भाग म्हणून येत असेल, तर पहिले पद दुसऱ्याचे उपपद आहे, अशी कल्पना करा. त्यासाठी $t_1 \sqsubseteq t_2$ हे चिन्ह वापरू. अर्थात, प्रत्येक नीटस पद स्वतःचे उपपद असते: $t \sqsubseteq t$. या संबंधाची विगमनाधारित व्याख्या अशी देता येते:

\begin{defn}
नीटस पद~$t_1$ हे~$t_2$ चे उपपद आहे हा संबंध, $t_1 \sqsubseteq t_2$, पुढीलप्रमाणे~$t_2$ वर विगमनाने परिभाषित केला आहे:
\begin{enumerate}
\item $t_2$ हे अक्षर असेल, तर $t_1 \sqsubseteq t_2$ हे $t_1 = t_2$ असेल तेव्हाच आणि फक्त तेव्हाच सत्य असते.
\item $t_2$ हे $[s_1 \circ s_2]$ असेल, तर $t_1 \sqsubseteq t_2$ हे $t_1 = t_2$, $t_1 \sqsubseteq s_1$ किंवा $t_1 \sqsubseteq s_2$ यांपैकी एक असेल तेव्हाच आणि फक्त तेव्हाच सत्य असते.
\end{enumerate}
\end{defn}

उदाहरणार्थ, या व्याख्येनुसार $\mathrm{a} \sqsubseteq [\mathrm{b} \circ \mathrm{a}]$ मिळते. कारण (2) नुसार $\mathrm{a} \sqsubseteq [\mathrm{b} \circ \mathrm{a}]$ हे $\mathrm{a} = [\mathrm{b} \circ \mathrm{a}]$, किंवा $\mathrm{a} \sqsubseteq b$, किंवा $\mathrm{a} \sqsubseteq \mathrm{a}$ यांपैकी एक असेल तेव्हाच आणि फक्त तेव्हाच सत्य असते. पहिले दोन पर्याय असत्य आहेत: $\mathrm{a}$ हे स्पष्टपणे $[\mathrm{b} \circ \mathrm{a}]$ शी एकरूप नाही; आणि (1) नुसार $\mathrm{a} \sqsubseteq \mathrm{b}$ हे $\mathrm{a} = \mathrm{b}$ असेल तेव्हाच आणि फक्त तेव्हाच सत्य असते; तेही असत्य आहे. पण पुन्हा (1) नुसार $\mathrm{a} \sqsubseteq \mathrm{a}$ हे $\mathrm{a} = \mathrm{a}$ असेल तेव्हाच आणि फक्त तेव्हाच सत्य असते; ते सत्य आहे.

ही व्याख्या यशस्वी होण्यासाठी एक अजून न सिद्ध केलेले तथ्य आवश्यक आहे: प्रत्येक नीटस पद~$t$ हे स्वतंत्र अक्षर असते, किंवा अशी \emph{एकमेव रीतीने निश्चित} होणारी नीटस पदे $s_1$ आणि $s_2$ असतात की $t = [s_1 \circ s_2]$. येथे “एकमेव रीतीने निश्चित” याचा अर्थ असा: $t = [s_1 \circ s_2]$ असेल, तर $s_1 \neq r_1$ किंवा $s_2 \neq r_2$ असताना ते \emph{त्याचबरोबर} $= [r_1 \circ r_2]$ असू शकत नाही. असे घडले, तर खंड~(2) चाच स्वतःशी संघर्ष होऊ शकतो: $t_2$ ला $[s_1 \circ s_2]$ असे वाचल्यास $t_1 \sqsubseteq t_2$ मिळेल; पण $t_2$ ला $[r_1 \circ r_2]$ असे वाचल्यास $t_1 \sqsubseteq t_2$ असत्य मिळेल. हे घडू शकत नाही हे सिद्ध करण्याआधी, जिथे ते \emph{घडू शकते} असे एक उदाहरण पाहू.

\begin{defn}
\emph{कंसविरहित पदे} विगमनाधारित रीतीने अशी परिभाषित करा:
\begin{enumerate}
\item प्रत्येक अक्षर कंसविरहित पद आहे.
\item $s_1$ आणि $s_2$ कंसविरहित पदे असतील, तर $s_1 \circ s_2$ हे कंसविरहित पद आहे.
\item इतर काहीही कंसविरहित पद नाही.
\end{enumerate}
\end{defn}

कंसविरहित पदांची उदाहरणे म्हणजे $\mathrm{a}$, $\mathrm{b} \circ \mathrm{d}$ आणि $\mathrm{b} \circ \mathrm{a} \circ \mathrm{b}$. यांच्यासाठी “उपपद” ची व्याख्या वरच्या पद्धतीने केली, तर दुसरा खंड असा असेल:
\begin{center}
$t_2 = s_1 \circ s_2$ असेल, तर $t_1 \sqsubseteq t_2$ हे $t_1 = t_2$, $t_1 \sqsubseteq s_1$ किंवा $t_1 \sqsubseteq s_2$ असेल तेव्हाच आणि फक्त तेव्हाच सत्य असते.
\end{center}

आता $\mathrm{b} \circ \mathrm{a} \circ \mathrm{b}$ हे $s_1 \circ s_2$ या रूपाचे आहे, जिथे
\begin{align*}
s_1 & = \mathrm{b} \text{ आणि} & s_2 & = \mathrm{a} \circ \mathrm{b}.
\intertext{तेच पद $r_1 \circ r_2$ या रूपाचेही आहे, जिथे}
r_1 & = \mathrm{b} \circ \mathrm{a} \text{ आणि} & r_2 &= \mathrm{b}.
\end{align*}
आता $\mathrm{a} \circ \mathrm{b}$ हे $\mathrm{b} \circ \mathrm{a} \circ \mathrm{b}$ चे उपपद आहे का? पहिल्या वाचनानुसार उत्तर होय; दुसऱ्यानुसार नाही.

एखादे नीटस पद इतर नीटस पदांपासून कसे रचले आहे हे एकमेव रीतीने निश्चित होते, या गुणधर्माला \emph{एकमेव रीतीने वाचनीयता} म्हणतात. अशा विगमनाधारित वस्तूंवरील संबंधांच्या विगमनाधारित व्याख्या महत्त्वाच्या असल्यामुळे हा गुणधर्म सिद्ध केला पाहिजे.

\begin{prop}
$t$ हे नीटस पद आहे असे समजा. मग $t$ हे स्वतंत्र अक्षर आहे; किंवा अशी एकमेव रीतीने निश्चित होणारी नीटस पदे $s_1$, $s_2$ आहेत की~$t = [s_1 \circ s_2]$.
\end{prop}

\begin{proof}
$t$ स्वतंत्र अक्षर असेल, तर अट पूर्ण होते. म्हणून $t$ स्वतंत्र अक्षर नाही असे समजा. विगमनाधारित व्याख्येवरून मग $t$ हे काही नीटस पदे $s_1$ आणि~$s_2$ यांसाठी $[s_1 \circ s_2]$ या रूपाचे असले पाहिजे. ही पदे एकमेव रीतीने निश्चित होतात हे दाखवायचे आहे: $t = [r_1 \circ r_2]$ असेल, तर $s_1 = r_1$ आणि $s_2 = r_2$.

म्हणून $t = [s_1 \circ s_2]$ आणि $t = [r_1 \circ r_2]$ असेही समजा; येथे $s_1$, $s_2$, $r_1$, $r_2$ ही नीटस पदे आहेत. $s_1 = r_1$ आणि $s_2 = r_2$ दाखवायचे आहे. प्रथम $s_1$ आणि $r_1$ एकरूप असले पाहिजेत; तसे नसतील, तर त्यांपैकी एक दुसऱ्याचा अरिक्त उचित पूर्वखंड असेल. पण \ref{mth:ind:sti:prop:initial} नुसार $s_1$ आणि~$r_1$ दोन्ही नीटस पदे असतील, तर हे अशक्य आहे. $s_1 = r_1$ असल्यावर $s_2 = r_2$ देखील स्पष्टपणे मिळते.
\end{proof}

फलनेही विगमनाधारित रीतीने परिभाषित करता येतात. उदाहरणार्थ, कोणत्याही नीटस पदातील आतल्या आत असलेल्या $[\dots]$ जोड्यांची कमाल खोली देणारे~$f$ फलन पुढीलप्रमाणे परिभाषित करता येते:

\begin{defn}
\label{mth:ind:rel:defn:depth} नीटस पदाची \emph{खोली}, $f(t)$, पुढीलप्रमाणे विगमनाधारित व्याख्येने दिली आहे:
\[
f(t) = \begin{cases}
0 & \text{ जर $t$ अक्षर असेल}\\
\max(f(s_1), f(s_2)) + 1 & \text{ जर $t = [s_1 \circ s_2]$ असेल.}
\end{cases}
\]
\end{defn}

उदाहरणार्थ,
\begin{align*}
f([\mathrm{a} \circ \mathrm{b}]) & =
\max(f(\mathrm{a}),f(\mathrm{b})) + 1 = \\
&= \max(0, 0) + 1 = 1, \text{ आणि}\\
f([[\mathrm{a} \circ \mathrm{b}] \circ \mathrm{c}]) & =
\max(f([\mathrm{a} \circ \mathrm{b}]), f(\mathrm{c})) + 1 = \\
& = \max(1,0) + 1 = 2.
\end{align*}

अर्थात, येथे $s_1$ आणि $s_2$ ही नीटस पदे आहेत असे मानतो. प्रत्येक नीटस पद स्वतंत्र अक्षर किंवा $[s_1 \circ s_2]$ या रूपाचे असते, हे तथ्य वापरतो. ते या रूपात एकाच रीतीने असू शकते, हेही पुन्हा महत्त्वाचे आहे. का ते पाहण्यासाठी आधी परिभाषित केलेल्या कंसविरहित पदांचा पुन्हा विचार करा. त्यांच्यासाठी यासारखी “व्याख्या” अशी असेल:
\[
g(t) =
\begin{cases}
0 & \text{ जर $t$ अक्षर असेल}\\
\max(g(s_1), g(s_2)) + 1 & \text{ जर $t = s_1 \circ s_2$ असेल.}
\end{cases}
\]
आता $\mathrm{a} \circ \mathrm{b} \circ \mathrm{c} \circ \mathrm{d}$ हे कंसविरहित पद पाहा. ते अनेक रीतींनी वाचता येते. उदाहरणार्थ, $s_1 \circ s_2$ या रूपात, जिथे
\begin{align*}
s_1 & = \mathrm{a} \text{ आणि} &
s_2 & = \mathrm{b} \circ \mathrm{c} \circ \mathrm{d},
\intertext{किंवा $r_1 \circ r_2$ या रूपात, जिथे}
r_1 & = \mathrm{a} \circ b \text{ आणि} &
r_2 &= \mathrm{c} \circ \mathrm{d}.
\end{align*}
पहिल्या वाचनानुसार $g$ मोजल्यास असे मिळेल:
\begin{align*}
g(s_1 \circ s_2) & =
\max(g(\mathrm{a}), g(\mathrm{b} \circ \mathrm{c} \circ \mathrm{d})) + 1 =\\ & =
\max(0,2) + 1 = 3
\intertext{तर दुसऱ्या वाचनानुसार असे मिळते}
g(r_1 \circ r_2) & =
\max(g(\mathrm{a} \circ \mathrm{b}), g(\mathrm{c} \circ \mathrm{d})) + 1=\\ &
= \max(1,1) + 1 = 2
\end{align*}
पण फलनाने नेहमी एकमेव मूल्य दिले पाहिजे. म्हणून~$g$ ची ही “व्याख्या” मुळात फलनच परिभाषित करत नाही.

\begin{prob}
$l$ फलनाची विगमनाधारित व्याख्या द्या; येथे $l(t)$ म्हणजे नीटस पद~$t$ मधील चिन्हांची संख्या.
\end{prob}

\begin{prob}
नीटस पद~$t$ वर संरचनात्मक विगमन वापरून $f(t) < l(t)$ सिद्ध करा. येथे $l(t)$ म्हणजे~$t$ मधील चिन्हांची संख्या आणि $f(t)$ म्हणजे \ref{mth:ind:rel:defn:depth} मध्ये परिभाषित केलेली~$t$ ची खोली.
\end{prob}



\part{सिद्धता उपपत्ती}\label{pt:part}
\begin{editorial}
  सिद्धता उपपत्तीवरील साहित्य अपूर्ण आणि प्रायोगिक आहे;
  ते अद्याप मुख्य PDF संकलनात समाविष्ट केलेले नाही.
\end{editorial}
